\section{Half-rectangle bijections}\label{sec:half}
We now derive the fixed-point counts without using the block
classification or the palindrome argument. The local construction
is stated for any positive height.

\subsection{Partial tilings and a transfer matrix}
For $m,j\ge1$ and $S\subseteq[m]=\{1,\ldots,m\}$,
let $\mathcal P_m(j,S)$ be the set of partial tilings of an
$m\times j$ rectangle that cover every cell except the cells in
rows $S$ of the last column.
Each uncovered cell is regarded as the left half of a horizontal
domino to be completed in the next column.
For $j=0$, only the empty object with $S=\varnothing$ is allowed.
Write $q_j(S)=|\mathcal P_m(j,S)|$.

Index rows and columns of $T_m$ by all subsets of $[m]$, and set
\begin{equation}\label{eq:Tm}
 T_m[p,q]=
 \begin{cases}
 1,&p\cap q=\varnothing\ \text{and every consecutive run}\\
   &\quad\text{in }[m]\setminus(p\cup q)\text{ has even length},\\
 0,&\text{otherwise}.
 \end{cases}
\end{equation}
An empty complement satisfies the condition.
The cells in $p$ complete the incoming horizontal dominoes;
those in $q$ remain as outgoing halves.
Each remaining consecutive run has a vertical tiling if and
only if its length is even, and then that tiling is unique.
This follows by pairing the endpoint with its only available
neighbour and proceeding along the run.
Consequently, legal state paths and partial tilings determine
one another uniquely.

Let $\evec$ be the unit column vector at the empty state, and
let $t_m(j)$ count complete $m\times j$ tilings. Then
\begin{equation}\label{eq:profile}
 q_j^{\trans}=\evec^{\trans}T_m^j,\qquad
 t_m(j)=\evec^{\trans}T_m^j\evec.
\end{equation}
For an auxiliary height-zero term, set $T_0=(1)$ and $t_0(j)=1$.

\subsection{Boundary conditions imposed by symmetry}
Define $J(S)=\{m+1-r:r\in S\}$.
Let $\one$ be the all-ones column vector, and let
$\chi_m(S)=1$ if $J(S)=S$, and $0$ otherwise.
Denote the fixed-point counts of the two reflections and the
half-turn on an $m\times2n$ rectangle by $h_{m,n},v_{m,n},r_{m,n}$,
respectively.

\begin{lemma}\label{lem:half}
For every $m\ge1$ and $n\ge0$,
\begin{equation}\label{eq:half-fixed}
 v_{m,n}=\evec^{\trans}T_m^n\one,\qquad
 r_{m,n}=\evec^{\trans}T_m^n\chi_m.
\end{equation}
\end{lemma}
\begin{proof}
Suppose first that $n\ge1$. Given $P\in\mathcal P_m(n,S)$,
reflect its internal dominoes into the right half of the
rectangle, and insert a horizontal domino joining $(r,n)$
to $(r,n+1)$ for each $r\in S$.
The internal dominoes in the two halves are disjoint, and the
inserted dominoes cover precisely the remaining cells.
The resulting tiling is fixed by vertical-axis reflection.
Conversely, restrict any such fixed tiling to its left half
and cut its midline-crossing dominoes. This recovers $P,S$
uniquely, proving the first bijection and formula.

For the half-turn, rotate the internal dominoes of $P$ into
the right half. The unmatched rows on the right are $J(S)$.
A domino crossing the midline must join two cells in the same
row, so completion is possible if and only if $J(S)=S$.
When this condition holds, inserting one crossing domino in
each row of $S$ gives a half-turn-invariant tiling.
Restriction to the left half is again the inverse map.
Equivalently, necessity follows because the half-turn sends a
crossing domino in row $r$ to one in row $m+1-r$.
The empty-object convention gives both formulas for $n=0$.
\end{proof}

\begin{lemma}\label{lem:h-odd}
If $m=2k+1$ with $k\ge0$, then
\begin{equation}\label{eq:h-odd}
 h_{m,n}=t_k(2n)\qquad(n\ge0).
\end{equation}
\end{lemma}
\begin{proof}
In a tiling fixed by horizontal-axis reflection, a vertical
domino incident with the middle row would have a distinct
reflected image sharing its middle-row cell. Such a domino
is impossible.
The middle row must therefore have its unique horizontal
tiling. The top $k$ rows form a complete tiling, and their
reflection determines the bottom $k$ rows.
Conversely, every tiling of the top subrectangle gives a fixed
tiling in this way. Restriction is the inverse.
The conventions cover $k=0$ and $n=0$.
\end{proof}

\subsection{A two-dimensional recurrence for height three}
From \eqref{eq:Tm}, the states reachable from the empty state
at height three are
\[
\varnothing,\quad\{1\},\quad\{3\},\quad\{1,2\},\quad\{2,3\},\quad\{1,2,3\}.
\]
The empty state leads to $\{1\},\{3\},\{1,2,3\}$;
$\{1\}$ leads to $\varnothing,\{2,3\}$;
$\{3\}$ leads to $\varnothing,\{1,2\}$;
The remaining transitions are
\[
\{1,2\}\longrightarrow\{3\},\qquad
\{2,3\}\longrightarrow\{1\},\qquad
\{1,2,3\}\longrightarrow\varnothing.
\]
We now make the induction explicit. Put
\[
 E=\{\varnothing,\{1,2\},\{2,3\}\},\qquad
 O=\{\{1\},\{3\},\{1,2,3\}\}.
\]
Every transition out of $E$ ends in $O$, and every transition out of
$O$ ends in $E$. There is no exit to either of the other two states.
This is a statement about the complete list of allowed exits, not a
claim that every pair of states in the two sets is connected.

\paragraph{Induction statement and base case.}
For each $j\ge0$, the support of $q_{2j}$ is contained in $E$, and
\[
 q_{2j}(\varnothing)=\alpha_j,\qquad
 q_{2j}(\{1,2\})=q_{2j}(\{2,3\})=\beta_j.
\]
Thus $\beta_j$ counts one specified state; the two nonempty states
together contribute $2\beta_j$. We do not assert that all states of
$E$ have positive counts. At width zero only the empty state occurs,
so $(\alpha_0,\beta_0)=(1,0)$, including equality of the two zero counts.

\paragraph{One added column.}
Assume the statement at width $2j$. Grouping by the previous state gives
\[
\begin{aligned}
q_{2j+1}(\{1\})
 &=q_{2j}(\varnothing)+q_{2j}(\{2,3\})=\alpha_j+\beta_j,\\
q_{2j+1}(\{3\})
 &=q_{2j}(\varnothing)+q_{2j}(\{1,2\})=\alpha_j+\beta_j,\\
q_{2j+1}(\{1,2,3\})&=q_{2j}(\varnothing)=\alpha_j.
\end{aligned}
\]
All other counts are zero because the exits from $E$ have been exhausted.
Each legal transition has exactly one local filling. Conversely,
removing the added column, and cutting any horizontal domino crossing
its left boundary, recovers a unique previous object and state.
The source classes are therefore disjoint and exhaustive; no extra
multiplicity is needed in these sums.

\paragraph{A second added column.}
All three states in $O$ can end at the empty state. Only $\{3\}$
can end at $\{1,2\}$, and only $\{1\}$ can end at $\{2,3\}$.
Consequently,
\[
\begin{aligned}
q_{2j+2}(\varnothing)
 &=q_{2j+1}(\{1\})+q_{2j+1}(\{3\})+q_{2j+1}(\{1,2,3\})\\
 &=(\alpha_j+\beta_j)+(\alpha_j+\beta_j)+\alpha_j
   =3\alpha_j+2\beta_j,\\
q_{2j+2}(\{1,2\})&=q_{2j+1}(\{3\})=\alpha_j+\beta_j,\\
q_{2j+2}(\{2,3\})&=q_{2j+1}(\{1\})=\alpha_j+\beta_j.
\end{aligned}
\]
There are no other nonzero counts, since all exits from $O$ lie in $E$.
This proves the support statement and equality of the two counts at
the next even width, completing the induction for every $j\ge0$.
In column-vector form the recurrence is
\begin{equation}\label{eq:two-dimensional}
 \binom{\alpha_{j+1}}{\beta_{j+1}}
 =\begin{pmatrix}3&2\\1&1\end{pmatrix}
 \binom{\alpha_j}{\beta_j},\qquad
 (\alpha_0,\beta_0)=(1,0).
\end{equation}

\paragraph{Reading the four coefficients.}
The first row counts arrivals at $\varnothing$; the second counts
arrivals at the single target $\{1,2\}$.
\begin{center}
\small
\begin{tabular}{cp{0.78\linewidth}}
\toprule
Entry & Two-column continuation counted\\
\midrule
$3$ & From $\varnothing$ back to $\varnothing$, through any of
$\{1\}$, $\{3\}$, or $\{1,2,3\}$.\\
$2$ & One return to $\varnothing$ from each of $\{1,2\}$ and
$\{2,3\}$, giving two source classes of size $\beta_j$.\\
$1$ & From $\varnothing$ through $\{3\}$ to $\{1,2\}$.\\
$1$ & From $\{1,2\}$ through $\{3\}$ back to $\{1,2\}$;
the other nonempty state contributes nothing to this target.\\
\bottomrule
\end{tabular}
\end{center}
In particular, the coefficient $2$ does not mean that each object
in a fixed nonempty state has two returns to the empty state.
For example, $(\alpha_1,\beta_1)=(3,1)$ gives
$\alpha_2=(3+1)+(3+1)+3=11$ and $\beta_2=3+1=4$.
This example illustrates the induction; it is not its basis.

\begin{proposition}\label{prop:parity}
The ordinary counts $d_j=t_3(2j)$ satisfy
\begin{equation}\label{eq:d-rec}
 d_0=1,\qquad d_1=3,\qquad
 d_{j+2}=4d_{j+1}-d_j\quad(j\ge0).
\end{equation}
For every $j\ge0$,
\begin{equation}\label{eq:parity-fixed}
\begin{gathered}
 h_n=1\quad(n\ge0),\qquad r_{2j}=r_{2j+1}=d_j,\\
 v_{2j}=d_{j+1}-2d_j,\qquad v_{2j+1}=d_{j+1}.
\end{gathered}
\end{equation}
In particular,
\begin{equation}\label{eq:parity-orbits}
\begin{split}
 a(2j)&=\frac{d_{2j}+1+d_{j+1}-d_j}{4},\\
 a(2j+1)&=\frac{d_{2j+1}+1+d_{j+1}+d_j}{4}.
\end{split}
\end{equation}
\end{proposition}
\begin{proof}
Equation \eqref{eq:profile} gives $d_j=\alpha_j$.
The matrix in \eqref{eq:two-dimensional} satisfies
$M^2-4M+I=0$, proving \eqref{eq:d-rec}.
At an even step, the only reachable $J$-invariant state is
$\varnothing$; at an odd step it is $\{1,2,3\}$.
Lemma \ref{lem:half} therefore gives the two formulas for $r$.
Summing all terminal counts gives
$\alpha_j+2\beta_j=d_{j+1}-2d_j$ at an even step and
$3\alpha_j+2\beta_j=d_{j+1}$ at an odd step.
Lemma \ref{lem:h-odd} gives $h_n=t_1(2n)=1$.
Substitution in Burnside's formula proves \eqref{eq:parity-orbits}.
\end{proof}

\begin{proof}[Second proof of Theorem \ref{thm:main}]
The recurrence \eqref{eq:d-rec} gives
$D(x)=(1-x)/(1-4x+x^2)$.
Summing \eqref{eq:parity-fixed} over the two parities gives
\[
 H(x)=\frac1{1-x},\qquad R(x)=(1+x)D(x^2),
\]
and
\[
 V(x)=(1+x)\frac{D(x^2)-1}{x^2}-2D(x^2)
 =\frac{1+3x+x^2-x^3}{1-4x^2+x^4}.
\]
The quotient denotes the unique coefficient shift: $D(t)-1=tE(t)$, where $E(t)=\sum_{j\ge0}d_{j+1}t^j$. It does not treat $t$ as a unit in the formal power-series ring.
Burnside's formula now gives \eqref{eq:average-rational};
the algebraic simplification and coefficient comparison in
the first proof finish the argument.
The combinatorial derivation here uses only the partial-tiling
model and the half-rectangle bijections.
\end{proof}
